\bookchapter{Bab 58: Manual rujukan pantas dan sumber}

\chapterintro{Bab ini ialah manual rujukan pantas yang boleh digunakan secara berasingan. Apabila anda lupa sesuatu sintaks atau perlu membuat keputusan pantas, terus rujuk sini. Disyorkan untuk menjadikannya “kad saku”: rujukan pantas kod biasa, soalan lazim, sumber pembelajaran — tiada halaman tertinggal.}

\booksection{Kad rujukan pantas kod biasa}

\begin{center}
\begin{tabularx}{\linewidth}{|X|X|}
\hline
\textbf{Apa yang mahu dibuat} & \textbf{Kod utama} \tabularnewline
\hline
Ingat satu keadaan boleh ubah & var x by remember \{ mutableStateOf(nilaiAwal) \} \tabularnewline
Keadaan tidak hilang apabila skrin diputar & rememberSaveable \{ mutableStateOf(nilaiAwal) \} \tabularnewline
Susunan menegak/melintang & \texttt{Column \{ \}}/\texttt{Row \{ \}} \tabularnewline
Senarai muat lambat & LazyColumn \{ items(senarai) \{ ... \} \} \tabularnewline
Navigasi antara halaman & navController.navigate("laluan") \tabularnewline
Kerja menunggu dilakukan dalam coroutine & \texttt{viewModelScope.\allowbreak launch \{ .\allowbreak .\allowbreak .\allowbreak  \}} \tabularnewline
Memerhatikan aliran data & \texttt{val list by flow.\allowbreak collectAsState()} \tabularnewline
Simpan konfigurasi kecil & context.dataStore.edit \{ it[KEY] = nilai \} \tabularnewline
Simpan data berstruktur & Tiga serangkai Room: @Entity / @Dao / @Database \tabularnewline
Ambil data melalui rangkaian & Antara muka Retrofit + fungsi suspend \tabularnewline
Hantar pemberitahuan & NotificationChannel + NotificationCompat.Builder \tabularnewline
Tulis log & \texttt{Log.\allowbreak d("TAG", "message")} \tabularnewline
\hline
\end{tabularx}
\end{center}

\begin{center}\small\color{gray}Kad rujukan pantas pembangunan Android (edisi Ogos 2026)\end{center}

\booksection{Soalan lazim (FAQ)}

S1: Saya langsung tiada asas pengaturcaraan, berapa lama untuk belajar pembangunan Android?

Ikuti laluan buku ini: minggu 1-2 pasang persekitaran dan belajar asas Kotlin; minggu 3-4 belajar Compose untuk membina antara muka; minggu 5-6 siapkan projek buku perakaunan; minggu 7-8 habiskan projek cuaca. Dua bulan, daripada sifar kepada mampu membuat aplikasi kecil secara berasingan. Syaratnya: dedikasikan 1-2 jam sehari, dan taip kod dengan tangan, bukan sekadar membaca.

S2: Adakah Java masih perlu dipelajari?

Tidak perlu. Projek baharu, tutorial baharu dan jawatan baharu pada asasnya semuanya menggunakan Kotlin. Anda hanya akan bertemu banyak Java apabila “menyelenggara projek lama berumur sepuluh tahun”. Selepas menguasai Kotlin, jika kerja memerlukannya, menambah Java akan cepat — kedua-duanya berkongsi ekosistem mesin maya yang sama.

S3: Emulator atau peranti sebenar?

Emulator memang mudah untuk debug semasa pembangunan (snapshot, pelbagai model), tetapi penerimaan akhir mesti melalui peranti sebenar — sentuhan, prestasi, rangkaian sebenar, penderia, semuanya tidak boleh digantikan oleh emulator. Jika berkemampuan, membeli satu peranti sebenar kelas pertengahan sebagai peranti ujian ialah pelaburan paling berbaloi.

S4: Senangkah mendapat kerja dalam pembangunan Android?

Keadaan pasaran 2026: jawatan peringkat permulaan yang sekadar “boleh tulis antara muka” sangat kompetitif, tetapi pembangun yang “faham Kotlin + mahir seni bina + ada karya yang telah diterbitkan” masih dicari-cari. Untuk anda yang masih belajar, resume terbaik ialah aplikasi buatan sendiri yang ada di kedai.

S5: Apa yang perlu dibuat jika terkandas kerana ralat?

Tiga langkah emas: (1) salin baris pertama ralat dan cari seadanya di StackOverflow/mesin carian; (2) semak halaman berkaitan dalam dokumentasi rasmi; (3) jika masih gagal, bungkus “tingkah laku diharapkan + gejala sebenar + teks penuh ralat + kod berkaitan” dan hantar kepada AI atau komuniti untuk bertanya. Sebahagian besar masalah selesai pada langkah pertama.

\booksection{Sumber pembelajaran disyorkan}

\begin{center}
\begin{tabularx}{\linewidth}{|X|X|X|}
\hline
\textbf{Jenis} & \textbf{Sumber} & \textbf{Sebab disyorkan} \tabularnewline
\hline
Rasmi & developer.android.com & Dokumentasi rasmi Android, semua tutorial Compose/Kotlin ada di dalamnya, rujukan paling sahih \tabularnewline
Rasmi & Laman rasmi Kotlin (kotlinlang.org) & Dokumentasi rasmi bahasa + tapak latihan dalam talian \tabularnewline
Rasmi & Google Codelabs & Kursus latihan langkah demi langkah dengan bimbingan penuh, sangat mesra pemula \tabularnewline
Komuniti & Stack Overflow (tag android) & Stesen pertama meminta bantuan untuk ralat \tabularnewline
Komuniti & Juejin / CSDN (zon Android) & Banyak artikel praktikal berbahasa Cina; berhati-hati membezakan yang baharu dan lama (semak tarikh dahulu) \tabularnewline
Video & Saluran rasmi Google Developers & Penerangan teknologi baharu daripada persidangan pembangun setiap tahun \tabularnewline
Temahan buku & Bahagian 7 "Amalan Pengaturcaraan AI" (Bab 37-47) & Panduan lengkap menggunakan AI untuk pembangunan Android \tabularnewline
\hline
\end{tabularx}
\end{center}

\begin{center}\small\color{gray}Senarai sumber pembelajaran pembangunan Android (Ogos 2026)\end{center}

\booksection{Rujukan pantas arahan Gradle biasa}

Biasanya klik butang sudah cukup, tetapi baris arahan kadangkala lebih cekap (terutamanya apabila digandingkan dengan alat AI atau talian paip CI):

\begin{center}
\begin{tabularx}{\linewidth}{|X|X|}
\hline
\textbf{Arahan} & \textbf{Fungsi} \tabularnewline
\hline
\texttt{.\allowbreak /\allowbreak gradlew assembleDebug} & Bina APK versi debug (pada Windows, buang ./) \tabularnewline
\texttt{.\allowbreak /\allowbreak gradlew assembleRelease} & Bina versi release (perlu persiapan tandatangan) \tabularnewline
\texttt{.\allowbreak /\allowbreak gradlew test} & Jalankan semua ujian unit \tabularnewline
\texttt{.\allowbreak /\allowbreak gradlew lint} & Semakan statik: cari bug berpotensi dan gaya kod yang tidak menepati piawaian \tabularnewline
\texttt{.\allowbreak /\allowbreak gradlew clean} & Kosongkan hasil binaan (versi kompilasi daripada "petua mulakan semula") \tabularnewline
\texttt{.\allowbreak /\allowbreak gradlew :app: dependencies} & Lihat pepohon kebergantungan (alat ampuh mengesan konflik versi) \tabularnewline
\hline
\end{tabularx}
\end{center}

\begin{center}\small\color{gray}Rujukan pantas arahan Gradle biasa\end{center}

\begin{tcbinfo}{Petua: jika binaan bermasalah, jalankan clean dahulu}
Peristiwa misteri “semalam boleh jalan, hari ini tiba-tiba tidak” — enam atau tujuh daripada sepuluh kali disebabkan cache binaan bercelaru. Jalankan ./gradlew clean kemudian bina semula; ini boleh menyembuhkan lebih separuh masalah “misteri”. Jika masih tidak boleh, mulakan semula Android Studio — petua ini paling berkesan digandingkan dengan File Invalidate Caches.

\end{tcbinfo}

\booksection{Penutup: daripada aplikasi pertama ke aplikasi ke-N}

\booksection{Penutup: daripada aplikasi pertama ke aplikasi ke-N}

Bila dilihat semula, anda daripada “tidak tahu apa itu Activity” kepada memasang persekitaran, menguasai Kotlin, melukis antara muka dengan Compose, menyambungkan aplikasi kepada pangkalan data dan rangkaian, kemudian menguji, membungkus dan menerbitkan — di laluan ini, anda telah menamatkan bahagian paling sukar daripada perjalanan sifar ke satu.

Dunia teknologi berubah sangat pantas: rangka kerja akan dinaik taraf, alat akan berganti, AI akan mengambil alih lebih banyak kerja “menaip papan kekunci”. Tetapi ada satu perkara yang tidak akan lapuk: kebolehan anda menguraikan idea yang kabar kepada langkah-langkah yang jelas, kemudian melaksanakannya langkah demi langkah. Alat ialah pengganda, pemikiran ialah asas. Pengganda boleh diganti, tetapi asas tidak boleh digoyang.

Anda bukan lagi sekadar "orang yang mahir guna telefon". Anda boleh mengubah idea dalam kepala menjadi aplikasi yang benar-benar berjalan pada ratusan juta peranti Android. Mulai hari ini, berkarya dengan papan kekunci—itulah romantisme pembangun Android.

\booksection{Rujukan}

\begin{enumerate}
  \item Android Developers. (2026). Android Developer Guides and Jetpack Compose Documentation. developer.android.com.
  \item Kotlin Foundation. (2026). Kotlin Language Documentation. kotlinlang.org.
  \item Android Open Source Project. (2026). Android Platform Version History and API Levels.
  \item Google Play. (2026). Play Console Help: Publishing and App Signing.
  \item AndroidX Release Notes. (2026). Room / DataStore / Navigation / WorkManager releases.
  \item Square. (2026). Retrofit Documentation. github.com/square/retrofit.
\end{enumerate}
